\chapter{Auction Theory}
\label{chap:auctions}

\section{Learning objectives}

After completing this chapter, you should be able to distinguish common auction formats, explain private and common values, identify the winner's curse, and assess basic auction-design objectives.

\section{Why auctions?}

Auctions allocate goods when prices are not known in advance or when buyers have different valuations. Governments use them for spectrum licences and procurement; firms use them for advertising and online sales; individuals use them for art, housing, and collectibles.

In a \keyterm{private-value} auction, each bidder knows their own valuation and it does not depend on others' valuations. In a \keyterm{common-value} auction, the item has approximately the same underlying value for all bidders, but bidders have different information about it. Oil leases and some artworks have important common-value elements.

\section{Auction formats}

An English auction has ascending public bids; the last bidder remaining wins and pays the final price. A Dutch auction descends from a high price until someone accepts. In a first-price sealed-bid auction, the highest bidder wins and pays their bid. In a second-price sealed-bid auction, the highest bidder wins but pays the second-highest bid.

In a private-value second-price auction, bidding one's true valuation is a dominant strategy: bidding higher risks paying more than the object is worth, while bidding lower risks losing an object that would have generated a surplus. The result depends on the format's assumptions and can fail with common values or additional strategic considerations.

\section{The winner's curse and design}

In common-value settings, winning may be bad news: the winner may have been the most optimistic bidder. This \keyterm{winner's curse} encourages rational bidders to shade bids below their estimates. Auction designers can reduce uncertainty through disclosure, inspection, reserve prices, and rules that discourage collusion.

Good auction design balances efficiency, revenue, simplicity, participation, and resistance to manipulation. These goals may conflict. A format that raises revenue in one setting need not allocate efficiently or attract enough credible bidders in another.

\section*{Chapter Review}

\subsection*{Key concepts}
\begin{itemize}
  \item \textbf{Private value}: bidder valuation known to that bidder and independent of others'.
  \item \textbf{Common value}: underlying value shared but uncertain.
  \item \textbf{Winner's curse}: tendency for a winner to overpay in common-value settings.
  \item \textbf{Second-price auction}: winner pays the second-highest bid.
\end{itemize}

\section*{Exercises}

\subsection*{Conceptual questions}
\begin{enumerate}[label=\arabic*.]
  \item Distinguish private-value from common-value auctions.
  \item Compare English, Dutch, first-price, and second-price auctions.
  \item Explain the winner's curse.
\end{enumerate}

\subsection*{Graphical exercises}
\begin{enumerate}[resume,label=\arabic*.]
  \item Sketch the price path in an English auction and a Dutch auction.
\end{enumerate}

\subsection*{Analytical exercises}
\begin{enumerate}[resume,label=\arabic*.]
  \item Explain why truthful bidding is optimal in a private-value second-price auction.
  \item A bidder estimates a common value at \euro100 but expects to be the most optimistic bidder if she wins. Explain why bidding \euro100 may be unwise.
\end{enumerate}

\subsection*{Exam-style multiple-choice questions}
\begin{enumerate}[resume,label=\arabic*.]
  \item In a second-price sealed-bid auction, the winner pays: (a) their own bid; (b) the reserve price only; (c) the second-highest bid; or (d) nothing.
  \item The winner's curse is especially relevant in: (a) common-value auctions; (b) private-value auctions with no uncertainty; (c) perfectly competitive markets only; or (d) public goods.
\end{enumerate}

\subsection*{PPE discussion question}

\noindent Should governments design auctions primarily to maximise public revenue or to allocate resources efficiently? Discuss using a public-spectrum or procurement example.
